\section{Introduzione all'IA}
\begin{quotation}
	Intendiamo l'Intelligenza Artificiale come l'abilità di saper risolvere dei problemi.
\end{quotation}
I due approcci per far risolvere a una macchina un problema sono dire esattamente la procedura risolutiva alla macchina o il \textbf{Machine Learning}, ossia la macchina impara da sola a risolvere i problemi dai dati.
\begin{center}
	\begin{minipage}{0.48\textwidth}
		\begin{adjustbox}{max width=\linewidth}
			\begin{tikzcd}
				\text{Inputs (observations)} \arrow[d] \arrow[rd] & {\text{[Approccio classico]}} \\
				Programmer \arrow[r] & Program \arrow[d] \\
				Outputs \arrow[u, bend left] & \langle \text{Computer} \rangle \arrow[l]
			\end{tikzcd}
		\end{adjustbox}
	\end{minipage}
	\hfill
	\begin{minipage}{0.48\textwidth}
		\begin{adjustbox}{max width=\linewidth}
			\begin{tikzcd}
				Inputs \arrow[rd] & {\text{[Machine Learning]}} & \\
				& \langle \text{Computer} \rangle \arrow[r] & Program \\
				Outputs \arrow[ru] & &
			\end{tikzcd}
		\end{adjustbox}
	\end{minipage}
\end{center}
Un sottoinsieme del Machine Learning è il Deep Learning che usa reti neurali per analizzare diversi fattori con una struttura simile al sistema neurale umano.
\subsection{Paradigma}
\begin{enumerate}
	\item DATA ACQUISITION, Elemento fondamentale e oggi obbiettivo centrale acquisirne in grandi quantità.
	\item DATA PROCESSING, Tecniche di elaborazione dati per adattarli al modello da sviluppare.
	\item MODEL, Nucleo di tecniche matematiche e statistiche in grado d'apprendere.
	\item PREDICTION, L'output del modello polimorfo a seconda del modello di cui va valutata efficacia tramite metriche appropriate.
\end{enumerate}
\paragraph{L'acquisizione di dati} avviene da risorse pubbliche o acquisendo un nuovo set di dati.

Trami l'annotazione o etichettatura si associa ai dati un'etichetta ossia la rappresentazione del contenuto semantico del dato e può essere numerica o categorica.
Un dato non annotato è inutile ameno che venga clusterizzato.
\begin{itemize}
	\item \textbf{Supervised learning} è un contestoin cui i dati sono annotati.
	\item \textbf{Unsupervised learning} è un contesto in cui i dati non sono annotati.

	L'algoritmo riconosce pattern nell'input senza indicazioni specifiche sugli output con lo scopo di autoapprendere la struttura degli input.

	Usato coi problemi di clustering per trovare gruppi di dati secondo caratteristiche comuni.
	\item \textbf{Semi-supervised learning} è un contesto in cui i dati sono miscelati. Non è sempre praticabile.
\end{itemize}
L'organizzazione dei dati è suddivisa in Training set, Validation set (dati su cui vengono messi a punto gli iperparametri) e Testing set.

\paragraph{Task} in ML a seconda dell'output atteso:
\begin{itemize}
	\item \textbf{Classificazione}, definisce la classe d'appartenenza di un input specifico. Se ci sono due classi si parla di classificazione binaria altrimenti di classificazione multiclasse.

	Una classe è un set di dati con proprietà comuni (semanticamente parlando) e è correlata all'etichetta.
	\item \textbf{Regressione}, trova relazioni tra le variabili dipendenti e indipendenti al fine di fare previsioni su nuovi dati.

	Restituisce un valore continuo.
	\item \textbf{Clustering}, identifica gruppi di input con caratteristiche simili.

	Usato in contesti non supervisionati. Il numero di gruppi non è noto a priori e spesso i gruppi diventano classi. Più complesso della classificazione.
\end{itemize}
\subsection{MLP}
Una rete neurale Multilayer Perceptron (MLP) è un insieme di perceptron, reti neurali di singoli neuroni artificiali, suddivise su più strati interconnessi tra loro. Un MLP può gestire separazioni più complicate di una lineare.

Una rete ha un livello di input, uno di output e diversi livelli intermedi detti hidden layer. Ogni livello è composto da diversi neuroni.
\paragraph{Ground-Trougt} è la veridicità della rete.
\subsubsection{Neurone Artificiale}
Il neurone artificiale imita l'architettura di un neurone biologico.
\begin{center}
	\begin{adjustbox}{max width=\textwidth}
		\begin{tikzcd}
			dendriti/input \arrow[rr] &  & corpo \ cellulare/somma \arrow[rr, "assone"] &  & sinapsi/output
		\end{tikzcd}
	\end{adjustbox}
\end{center}
\begin{itemize}
	\item \textbf{Dentriti}, dette estensioni cellulari, ricevono segnali elettrici/chimici da altri neuroni.
	\item \textbf{Il corpo} elabora i segnali dei dentriti nel tempo e converte il valore elaborato in un segnale.

	L'elaborazione consiste in una funzione di \textbf{soglia} che somma tutti gli input e se questi raggiungono una certa soglia in una finestra temporale, viene mandato un outut, altrimenti nulla.
	\item \textbf{L'assone} è il prolungamento del neurone hai neuroni, anche distanti, a cui trasmettere il segnale.
	\item \textbf{Le sinapsi} sono le estemità trasmissive dell'assone connesse hai neuroni, tramite i rispettivi dentriti.
\end{itemize}
Le sinapsi possono modificare individualmente la propria risposta modificando \textbf{l'efficienza di trasposto} dell'informazione.

Il valore di soglia e l'efficenza di trasporto per uno stesso neurone sono variabili e legate al processo d'apprendimento.
\subsubsection{Perceptron}
Il percettrone è il primo esempio di rete neurale. La rete è composta da un'unico neurone artificiale sul cui output viene applicata la funzione d'attivazione.
Le reti neurali normlai hanno diversi neuroni con il seguente.
\begin{center}
	\begin{adjustbox}{max width=\textwidth}
		\begin{tikzcd}
			in_1 \arrow[rd, "w_{1i}" description] &                                                      & out_i=f(net_i)                                                          \\
			\vdots \arrow[r, "\dots" description] & \sum \arrow[r]                                       & {net_i=\sum_{j=1,\dots,d}w_{ji}\cdot in_j+w_{0i}} \arrow[u, bend right] \\
			in_d \arrow[ru, "w_{di}" description] & in_0=1\text{ (bias)} \arrow[u, "w_{0i}" description] &
		\end{tikzcd}
	\end{adjustbox}
\end{center}
\paragraph{Funzinamento} Il neuore restituisce output $out_i$ della \textbf{funzione d'attivazione} $out_i=f(net_i)$ applicata sulla somma $net_i$ dei vari input e del \textbf{bias} moltiplicati per il rispettivo peso \textbf{peso}.

NOTA: La somma $net_i$ può essere sia positiva che negativa.
\subparagraph{Il peso} è l'importanza, quanto "accellero" o "rallento", che il neurone da a quell'input specifico. La $i$ indica il peso per l'i-esimo neurone.

Durante la fase d'adddestramento/apprendimento della rete si calibrano i pesi.
\subparagraph{Il bias} è il peso di un input fittizzio fissato a 1 di default. Serve a tara il punto ottimale di lavoro ponendo un ofset, addestrabile. In una rete moderna è opzionale e viene usato solamente nei neuroni teminali della rete.
\subparagraph{La funzione d'attivazione} determina se un neurone deve attivarsi, se input superiori a una certa soglia, e restituire un segnale o meno.
Le funzioni devono essere \textbf{non lineari} e \textbf{derivabili}.

Alcune funzioni: Sigmoid ($\frac{1}{1+e^{-x}}$), tanh ($\tanh(x)$), ReLu ($max(0, x)$), leaky ReLu ($max(0.1x, x)$), Maxout, ELU.
\subsubsection{Connessioni}
Nuove sinapsi/collegamenti vengono constantemente create in modo casuale.

Qunado due neuroni connessi si attivano simutaneamente il collegamento si rinforza e impariamo (aumentiamo il peso).

Se due neuroni connessi non si attivano mai assieme il collegamento si indebolisce e disimpariamo (diminuiamo il peso) fino a sciogliersi (le sinapsi si staccano). I due neuroni staccati possono riconettersi.

In base alle connessioni della rete possiamo avere:
\paragraph{FFNN} Le reti feedforward (allineamento in avanti) le connessioni collegano i neuroni di un livello coi neuroni di un livello successivo, mai di un livello precedente o lo stesso livello.
\paragraph{Ricorrenti} Le reti riccorrenti prevedono connessioni di feedback, verso neuroni dello stesso livello e a livelli precedenti.

I risultati dipendono dall'input e da una memoria (ingressi precedenti) della rete.
\subsubsection{Apprendimento}
Fissata la topologia della rete neurale l'addestramento è trovare i pesi che permettono il \textbf{mapping desiderato} tra input e output.
Per apprendere dai dati bisogna seguire le seguenti fasi:
\begin{center}
	\begin{adjustbox}{max width=\textwidth}
		\begin{tikzcd}
			apprendo \arrow[r] & provo \arrow[r, bend right] \arrow[rr, bend right] & correggo \arrow[l, bend right] & fine
		\end{tikzcd}
	\end{adjustbox}
\end{center}
Provare e correggere si chiama \textbf{apprendimento iterativo}. Apprendere piccole parti per volta si chiama \textbf{apprendimento per batch}.

Nota: una rete viene izializzata con pesi inizializzati in un intorno di 0, sia positivi che negativi.
\subsubsection{Funzione di loss}
La funzione di loss è usata nella fase dell'apprendimento per corregge la rete neurale prima della prossima ripetizione. \textbf{Misura l'errore della previsione del modello rispetto ai dati di training}. La funzione di loss deve essere derivabile.

Usasta per capire in automatico se sto incrementando (mi allontano) o diminuendo (mi avvicino) l'errore, tramite la derivata della funzione obiettivo. Apprendere significa minimizzare la funzione obiettivo.
\paragraph{La Discesa del gradiente} è l'avvicinarsi all'obiettivo/minimo a ogni "passo". Bisogna trovare la giusta dimensione del passo:
\begin{itemize}
	\item \textbf{Passi lunghi} possono saltare l'obiettivo.
	\item \textbf{Passi corti} richiedono troppo tempo per raggiungere l'obiettivo.
\end{itemize}
Utilizza Backpropagation e la regola della catena. Va ad'ottimizzare il \textbf{Learning rate} e il \textbf{Batch size}.
\paragraph{Binary Cross-Entropy} esempio di funzione di loss usata per i problemi di classificazione binaria \[BCE=-y_i\cdot \log\hat{y}_i-(1-y_i)\cdot\log(1-\hat{y}_i)\] dove $y_i$ è il vlaore attesso e $\hat{y}_i$ è la previsione dell'i-esimo neurone.

Scopo: punire molto gli errori molto gravi e punire poco gli errori lievi.

Funzionamento: $BCE$ è un valore d'errore sempre positivo dato solo da uno dei due membri, uno dei due si annulla, e i meno fanno si che il risultato sia positivo. Sono usati dei $\log$ perchè ci interessano gli output negativi, che crescono tanto in valore assoluto, dati dagli input tra 0 e 1.
\paragraph{Categorical Cross-Entropy} esempio di funzione di loss usata per i problemi di cassificazione multi classe \[CCE=-\sum_iy_i\cdot \log\hat{y}_i\]
Scopo: dare un valore basso per classe corretta.
Funzionamento: tutti valori sommatoria si annullano tranne uno. In un One-hot vector viene messo a 1 l'unico bit corrispondente alla classe indovinata e tutti gli altri stanno a 0.
\paragraph{Mean Squared Error o MSE} esempio di funzione di loss usata per i task di regressione \[C(W)=\frac{\sum_{i=1}^n(y_i-\hat{y}_i(W))^2}{n}\] Definito come la media dei quadrati delle differenze tra l'output previsto dal modello e l'output reale associato ai dati di training.
\begin{itemize}
	\item n = numero di esempi di training.
	\item $y_i$ è l'output reale associato a ciascun esempio di training.
	\item $\hat{y}_i(W)$ è l'output previsto dal modello per l'input corrispondente; con $W$ sono indicati i parametri della rete.
\end{itemize}
\paragraph{Mean Absolute Error o MAE} esempio di funzione di loss usata in problemi di regressione \[C(W)=\frac{\sum_{i=1}^n|y_i-\hat{y}_i(W)|}{n}\]
\paragraph{Backpropagation} dice dove abbiamo sbagliato.
\paragraph{Regola della catena} ???
\subsubsection{Softmax}
La rete nel layer di output da valori positivi o negativi a "caso", ma per le funzioni di loss servono percentuali.

Si aggiungere il layer \textbf{softmax} che trasforma l'output della rete (logits) in probabilità (valori in [0, 1[  che sommati danno 1).

Softmax è una funzione continua differenziabile che richiede n.neuroni output uguale al n.classi \[p_i=\frac{e^{a_i}}{\sum_{k=1}^ne^{a_k}}\]

Nota: tramite l'elevamento esponenziale si trasportano gli output negativi in positivi. Può gestire sia output molto grandi che molto piccoli. Indipendente dall'ordine di grandezza degli output della rete.
\subsection{CNN}
In un normale MLP ogni neurone di un layer è connesso a tutti i neuroni del layer precedente e un approccio tale con l'elaborazione delle immagini è pesante.

Le MLP non hanno alcuna invarianza per traslazione, ossia del dato che viene codificato da un neurone si perde la relazione di vicinanza a un altro dato assegnato a un altro neurone.
\paragraph{Le Convolutional Neural Network (CNN)} sono reti progettate per processare immagini.

Organizzano i neuroni in tre dimensioni. Collegano i neuroni localizatamente, vengono legati i neuroni ai neuroni associati ai pixel vicini e non a tutti.
\paragraph{La struttura} delle CNN è 3D, larghezza (W), altezza (H) e profondità (C) maggiore di 1.
\subsubsection{Kernel e convoluzione}
\paragraph{Il kernel} è un quadrato ($3\times 3$ tipicamente) con dei pesi fissati preallenabili. Praticamente è usato come filtro digitale sulle immagini per estrarre feature.
\paragraph{La convoluzione} è il processo matematico che collega il kernel all'immagine. Tramite uno \textbf{sliding-window}, il kernel si centra su ogni pixel dell'immagine uno dopo l'altro.
Per i pixel del bordo viene applicato un padding a 0.

\paragraph{Il feature map} è l'unione dei prodotti scalari tra il kernel e le porzione di immagine coperate durante lo sliding-window \[g(x,y)=(f*h)(x,y)=\sum_{i=0}^{F-1}\sum_{j=0}^{F-1}f(x,y)h(x-i, y-j)\] condensa una caratteristica dell'immagine.

Nota: Una profondità maggiore di 1 significa usare più kernel e ottenere più feature map che definiscono il volume di output.
\paragraph{I layer} convolutive più si allontanano dall'input più estraggono informazioni complesse.
\subsubsection{Procedimento}
\begin{center}
	\begin{adjustbox}{max width=\textwidth}
	\begin{tikzcd}
			Input \arrow[r] & Feature extractor \arrow[r] & Flatten \ layer \arrow[r] & Classifier \arrow[r] & output
		\end{tikzcd}
	\end{adjustbox}
\end{center}
\paragraph{Il Feature extractor} è una CNN. Tramite vari \textbf{Polling layers} e \textbf{2D convolution layers} vengono prodotte varie feature map, per analizzare aspetti differenti dell'immagine di input.
\subparagraph{Polling layers} riducono la dimensione spaziale dell'immagine (la comprime). Per esempio prendende il valore massimo o medio di una finestra dell'immaggine.

Vantaggi: invarianza alla posizione, riduce computazioni, previene overfitting.
\subparagraph{2D convolution layers} applicano i kernel all'immagine per estrarne le feature.
\paragraph{Il flatten layer} è ottenuto srotolando i feature map prodotto dal Feature extractor in un vettore di input per il Classifier.
\paragraph{Il Classifier} è un MLP con neuroni fully connected. Lo strato di output viene poi processato da una softmax activation function.
\subsubsection{Tipologie di training}
\paragraph{From scratch} addestra la rete da una configurazione random dei pesi tramite molti dati di training. Lungo e complesso.
\paragraph{Pre-trained} usa una rete gia addestrata senza modificare l'output originale. Non serve training ne dati.
\paragraph{Fine-tuned} usa rete pre addestrata congelando i pesi del feature extractor e cambio solo i pesi del classifier.